\section{Justage des Versuchs und Durchführung}
\subsection{Durchführung}
Die Funktion des Aufbaus wird geprüft, um die Datennahme mit optimaler Auflösung zu gewährleisten. 
Zur ersten Signalprüfung wird das Spektrometer auf 0° ausgerichtet und das Rohsignal des Szintillationsdetektors auf einem Oszilloskop betrachtet. Dabei ist die $^{137}\text{Cs}$-Quelle freigelegt und das Target abgeklemmt. Um Reflexionen und die daraus resultierende Signalverzerrung zu vermeiden, wird ein 50-$\Omega$-Abschlusswiderstand vom Oszilloskop an das Koaxialkabel angeschlossen; alternativ kann ein Koaxial-T-Stück mit externem Abschlusswiderstand verwendet werden.

Anschließend wird die Spitzenspannung der Photomultiplierröhre (PMT) schrittweise erhöht. Die Elektronenlawinenverstärkung der PMT wird so eingestellt, dass die Eingangsfrequenz und die Impulsamplitude im linearen Betriebsbereich liegen. Aufgrund der Potenzialdifferenzen der Detektorkomponenten und der Verkabelung erfolgt die Spannungserhöhung schrittweise, um eine Überlastung der PMT und instabile Betriebsbedingungen zu vermeiden.

Anschließend wird die Hochspannung der Photomultiplierröhre (PMT) schrittweise erhöht. Zur weiteren Signalverarbeitung wird ein Hauptverstärker (unipolarer Ausgang) zwischen PMT und Oszilloskop angeschlossen. Die Grob- und Feinverstärkung sowie die Impulsformungszeit sind an diesem Verstärker einstellbar, wodurch Hintergrundrauschen, Impulsbreite und spektrale Auflösung des Systems präzise beeinflusst werden.

Das geformte und verstärkte Signal wird digitalisiert und mit einem Vielkanalanalysator (MCA) und der Software MCA3 spektral analysiert. Zur Bestimmung des geeigneten Messfensters wird die spezifische 661-keV-Gammalinie im Cäsiumspektrum verwendet; die Verstärkung wird so lange angepasst, bis diese Linie im gewünschten Kanalbereich erscheint.

Um die Emissionslinien der $^{152}\text{Eu}$-Probe und die 661-keV-Linie des $^{137}\text{Cs}$-Isotops gleichzeitig mit optimaler Auflösung aufzuzeichnen, wird der Energiebereich auf etwa 1100 keV eingestellt.

\subsection{Beobachtungen}